% =====================================================================
% Section IV — Experiments (body-only; \input into root.tex)
% All numbers sourced from notes/fact_experiments.md (2026-07-15).
% Internal version codes must never appear here.
% =====================================================================
\section{Experiments}
\label{sec:experiments}

\subsection{Setup}
\label{sec:setup}

\textbf{Dataset and metric.}
We evaluate on DELIVER~\cite{cmnext2023}, a synthetic driving benchmark
with four spatially aligned modalities (RGB, depth, event, LiDAR),
25 semantic classes, and 3{,}983 / 2{,}005 / 1{,}897 images for
training / validation / test at $1042\times1042$ resolution. The data
cover five environmental conditions (cloudy, foggy, night, rainy,
sunny) as well as partial sensor-failure corruption cases. We report
mean Intersection-over-Union (mIoU) on both the validation and test
splits, and use the four-modality (camera--LiDAR--depth--event)
protocol throughout. MUSES~\cite{muses2024} results are planned:
\todo{MUSES official-protocol numbers (test-server submission pending;
only a non-comparable internal letterbox evaluation exists)}.

\textbf{Implementation details.}
The backbone is a single frozen DINOv3 ViT-L/16~\cite{dinov3_2025}
shared by all modalities; only the per-modality LoRA adapters
(rank~8), the fusion stack, the auxiliary decoders, and the
segmentation head are trained (46.8\,M trainable of 349.9\,M total,
13.4\%; the router adds ${\approx}0.27$\,M, \todo{exact full-model
parameter count from its training log}). Inputs are $1024^2$ crops.
We train for 200 epochs with AdamW (learning rate $6\times10^{-4}$,
weight decay 0.01), a warmup-poly schedule (power 0.9, 10-epoch
warmup), batch size 4 per GPU on 4 GPUs (effective batch 16 with
gradient accumulation), and bfloat16 mixed precision. The loss is
OHEM cross-entropy~\cite{ohem2016} plus the auxiliary
($\lambda_{\mathrm{aux}}{=}0.5$), calibration
($\lambda_{\mathrm{cal}}{=}0.1$), and router
($\lambda_{\mathrm{reg}}{=}0.01$, full model only) terms of
Sec.~\ref{sec:method}. Inference is single-scale without flipping or
any test-time augmentation. All experiments use a single seed (3407);
we therefore refrain from building claims on sub-1-point differences.

\textbf{Checkpoint-selection protocol.}
Both splits are evaluated every two epochs during training, which
creates a subtle risk of implicit test-set peeking when the reported
checkpoint is chosen post hoc. We adopt a strict protocol: all
headline numbers use the checkpoint with the best \emph{validation}
mIoU, and any test-selected (oracle) number is explicitly labeled as
such and excluded from state-of-the-art claims. Under this protocol
the no-router variant scores 68.19 val / 56.64 test (epoch~120); its
oracle test-best checkpoint reaches 57.60 (epoch~140) and is disclosed
only for completeness. One asymmetry remains: the no-router variant
was trained with a physics-based augmentation (PhysAug) that
competitors do not use, whereas the full model was trained without it;
\todo{PhysAug-free no-router re-run (T1) for a fully harmonized
ablation pair}.

\subsection{Comparison with State of the Art}
\label{sec:sota}

Table~\ref{tab:sota} compares ReliaDINO with published DELIVER
results. Following the protocol audit of Sec.~\ref{sec:related}, we
strictly separate the official \emph{test}-split cluster from
validation-only results and never mix the two; the ``extra
supervision'' column lists everything each method requires beyond
segmentation labels.

\begin{table}[t]
\centering
\caption{Comparison on DELIVER (four modalities: camera, LiDAR, depth,
event; 25 classes). ``Extra sup.''\ = supervision or data beyond
segmentation labels. Val and test columns are the official splits;
``--'' = not reported. $\dagger$: oracle test-selected checkpoint,
disclosed but excluded from claims (Sec.~\ref{sec:setup}).}
\label{tab:sota}
\setlength{\tabcolsep}{2.5pt}
\scriptsize
\resizebox{\columnwidth}{!}{%
\begin{tabular}{@{}llccc@{}}
\toprule
Method & Backbone & Extra sup. & Val & Test \\
\midrule
CMNeXt~\cite{cmnext2023}          & MiT-B2       & none                & 66.3  & 53.0 \\
GeminiFusion~\cite{geminifusion2024} & MiT-B2    & none                & 66.9  & 54.5 \\
MAGIC~\cite{magic2024}            & SegFormer-B2 & none                & 67.66 & --   \\
StitchFusion~\cite{stitchfusion2024} & MiT-B2    & none                & 68.18 & 53.4 \\
OmniSegmentor~\cite{omnisegmentor2025} & DFormer-L & mm.\ pretraining  & 68.0  & --   \\
CAFuser-CA$^2$~\cite{cafuser2025} & Swin-T       & cond.\ text (CLIP)  & 67.8  & 55.6 \\
CAFuser-CAA~\cite{cafuser2025}    & Swin-T       & cond.\ text (CLIP)  & 68.6  & 55.2 \\
DGFusion~\cite{dgfusion2026}      & Swin-T       & depth GT + cond.\ text & \todo{val n/a} & 56.7 \\
\midrule
ReliaDINO (no router)             & frozen DINOv3-L & none & 68.19 & 56.64\,/\,57.60$^\dagger$ \\
\textbf{ReliaDINO (full)}         & frozen DINOv3-L & none & \todo{67.74} & \todo{$\geq$57.14} \\
\midrule
\multicolumn{5}{@{}l@{}}{\todo{MUSES block (planned): DGFusion 79.5, CAFuser-CAA 78.5, ReliaDINO t.b.d.}}\\
\bottomrule
\end{tabular}%
}
\end{table}

On the test split, the full model reaches
\todo{57.14 (best checkpoint at the time of writing; training
completes on schedule and the final val-selected pair will replace
this number)}, exceeding the strongest published test result,
DGFusion at 56.7~\cite{dgfusion2026}, and the no-router variant's
oracle checkpoint reaches 57.60. Crucially, this is achieved with
segmentation labels \emph{only}: DGFusion additionally consumes
ground-truth log-depth supervision and CLIP-encoded condition
meta-text, and CAFuser requires condition annotations with
verbo-visual contrastive training. On the validation split, the
no-router variant (68.19) is within 0.41 of the best published result
(CAFuser-CAA, 68.6) while surpassing it by $+1.4$ on test. We note the
concurrent MM SAM-adapter~\cite{mmsamadapter2025}, which reports
strong DELIVER numbers (e.g., 57.35 RGB-D) under a \emph{two-modality}
protocol with a fine-tuned side network; our claims are scoped to the
four-modality, supervision-minimal setting. A residual caveat: an
eval-time resize-convention parity check against DGFusion is
outstanding, \todo{confirm eval-protocol parity (G0d) before quoting
head-to-head margins below ${\sim}2$ points as definitive}.

\subsection{Ablation Study}
\label{sec:ablation}

Table~\ref{tab:ablation} groups the ablations into training-time
variants (top) and inference-time module toggles applied to a fixed
checkpoint (bottom).

\begin{table}[t]
\centering
\caption{Ablations on DELIVER. Top: trained variants (val-selected
unless noted; $\dagger$ = oracle test-selected). Bottom:
inference-time toggles on the no-router test-best checkpoint,
full-resolution full test set. $^a$trained with PhysAug (see
Sec.~\ref{sec:setup}); $^b$run terminated by a system event at epoch
120/200; $^c$attention biases were active during this variant's
training but are inert at inference (toggle $\Delta{\approx}0$).}
\label{tab:ablation}
\setlength{\tabcolsep}{3pt}
\scriptsize
\resizebox{\columnwidth}{!}{%
\begin{tabular}{@{}lcc@{}}
\toprule
Trained variant & Val & Test \\
\midrule
ReliaDINO (full: router + gate + calibration) & \todo{67.74} & \todo{57.14} \\
\;\;no router, harmonized recipe$^b$ & 67.61 & 56.14 \\
\;\;no router, original recipe$^{a,c}$ & 68.19 & 56.64\,/\,57.60$^\dagger$ \\
\midrule
Inference toggle (base: test 56.64) & \multicolumn{2}{c}{$\Delta$ test mIoU} \\
\midrule
$-$ attention-logit reliability bias ($\lambda_1 B^{\mathrm{cal}}$) & \multicolumn{2}{c}{$|\Delta|\leq0.03$ (dead)} \\
$-$ consistency bias ($\lambda_2 B^{\mathrm{cons}}$) & \multicolumn{2}{c}{$|\Delta|\leq0.03$ (dead)} \\
$-$ competence gate + calibration & \multicolumn{2}{c}{$-0.26$ (56.64$\rightarrow$56.38)} \\
$-$ veto floor & \multicolumn{2}{c}{$+0.02$\,--\,$+0.23$ (below claim thresh.)} \\
\bottomrule
\end{tabular}%
}
\end{table}

\textbf{Router.}
The per-class reliability-anchored router is the difference between
the full model and the harmonized-recipe no-router variant
(same recipe, incomplete run): $+0.13$ val / $+1.00$ test at the
respective best checkpoints so far
(\todo{final router pair after training completes}). The router's
benefit concentrates on the test split and appears early: at matched
epochs it led the no-router model by $+2.0$ to $+2.4$ test mIoU
(epochs 30--46) and crossed the DGFusion test mark at epoch 58 versus
epoch 116 without the router. The comparison against the
original-recipe no-router variant is confounded by its PhysAug
training (footnote $a$), which is why we report both rows.

\textbf{Competence gate and calibration.}
Disabling the calibrated competence gate at inference costs $-0.26$
test mIoU on the full test set at full resolution. The veto floor
contributes only within our single-seed noise band and we do not claim
it. The train-time calibration loss is not separable by an inference
toggle; we claim only its diagnostic effect (balanced per-modality
reliability AUROC, Sec.~\ref{sec:analysis}), not an mIoU gain.

\textbf{Attention-logit reliability biasing: a negative result.}
Across five model generations spanning two backbone families (SAM2
memory attention and the present cross-modal attention), injecting
reliability as a pre-softmax additive bias on attention logits never
helped: on the current architecture the toggle changes test mIoU by at
most $0.03$ with fused-feature cosine similarity of $1.0$ (the bias is
functionally inert), and on the strongest SAM2-generation variant it
was a statistically significant net \emph{harm}
($p=4.5\times10^{-22}$, paired per-image test). We report this
placement finding as a contribution rather than hiding it: once
per-modality LoRA absorbs modality adaptation, reliability information
is useful where class decisions are formed (output gate and per-class
router), not inside the attention operator. Removing the bias also
restores the fused attention kernel at inference.

\textbf{Frozen-backbone choice.}
A controlled linear-probe on frozen backbones (identical head and
data) gives DINOv3 ViT-L 40.51 test mIoU on RGB versus 28.90 for the
SAM2 Hiera-B+ encoder ($+11.6$), and 35.48 versus 25.34 on depth
($+10.1$); on thin structures the gap is starker (TrafficLight IoU
32.2 vs.\ 8.27). Consistently, our four generations of SAM2-based
reliability designs gained only $+0.7$ val mIoU in total, whereas the
backbone swap alone contributed $+4.1$: fused-feature effective rank
rises from 1.26 to 10.18 and cross-modal CKA from 0.02--0.16 to
0.80--0.91, i.e., the frozen VFM removes the feature collapse that
earlier reliability mechanisms were implicitly compensating for.

\textbf{Per-modality LoRA.}
Disabling a modality's LoRA adapters (frozen-backbone features
instead) drops that modality's auxiliary pixel accuracy in all five
conditions: by $+0.157$ to $+0.202$ for LiDAR and $+0.110$ to $+0.153$
for depth (the modalities farthest from the RGB pretraining
distribution), and by $+0.032$ to $+0.064$ (RGB) and $+0.020$ to
$+0.050$ (event). All 48 adapter pairs (24 blocks $\times$ 2
projections) remain active for every modality (0/48 dead), confirming
that the single frozen encoder is genuinely multiplexed rather than
defaulting to its RGB prior.

\subsection{Per-Class and Per-Condition Analysis}
\label{sec:analysis}

\textbf{Per-condition test audit.}
On the no-router test-best checkpoint, per-condition test mIoU is
55.64 (cloud), 56.32 (fog), 54.56 (night), 56.31 (rain), and 55.41
(sun) --- a spread of only 1.76 points. The large val--test gap of
DELIVER is therefore \emph{not} an adverse-condition effect but a
per-class transfer effect. To our knowledge this is the first
published per-class, per-condition audit of the DELIVER \emph{test}
split; competitors publish only aggregate test numbers, so
Table~\ref{tab:sota} has no external per-condition column.

\textbf{Benchmark ceiling.}
Three classes are effectively dead for \emph{all} methods we could
audit: Other (max IoU 5.6 across conditions), Wall (8.7), and Bridge
(0.5) score near zero for our model, for the frozen DINOv3 linear
probe, and for the official CMNeXt checkpoint
(\todo{exact official CMNeXt per-class test numbers from the M0-a
audit artifact}). We flag these as a benchmark ceiling: no fusion
mechanism can be credited or blamed for them. Conversely, classes that
were dead in our SAM2-generation models are recovered by the present
design: Water $0.0\rightarrow5.4$, TrafficLight $12.4\rightarrow36.6$,
Static $29.4\rightarrow39.1$ (five-condition means).

\textbf{Reliability quality.}
The calibration objective yields \emph{balanced} per-modality
reliability: at night, the AUROC of the calibrated confidence as a
correctness predictor is 0.85 / 0.78 / 0.87 / 0.70 for RGB / depth /
event / LiDAR. Earlier generations either collapsed on non-RGB
modalities (e.g., 0.84 / 0.63 / 0.26 / 0.38) or repaired non-RGB
geometry at the cost of RGB (0.43 / 0.92 / 0.55 / 0.97); ours is the
only generation with all four modalities above 0.70 without repair.

\textbf{Modality contributions.}
Drop-modality deltas (mIoU loss when a modality is removed at
inference) show depth as the dominant complement ($+9.5$ to $+11.4$
across conditions) and RGB rising from $+2.8$ (cloud) to $+8.0$
(night); LiDAR adds ${\leq}1.2$, and event contributes ${\approx}0$ in
all five conditions --- a known, lineage-wide limitation that fusion
alone does not fix. Qualitative results under adverse conditions are
shown in Fig.~\ref{fig:qual}, and the per-class audit, AUROC balance,
and gating behavior are visualized in Fig.~\ref{fig:analysis}.

\subsection{Fairness Protocol}
\label{sec:fairness}

For a clean reading of Table~\ref{tab:sota} we disclose all known
asymmetries. \emph{Ours:} a larger frozen backbone (349.9\,M total
parameters), though our 46.8\,M trainable parameters are fewer than
the ${\sim}58.7$\,M \emph{total} of CMNeXt-B2 and comparable to
trainable Swin-T pipelines (${\sim}79$\,M); the no-router variant used
PhysAug (Sec.~\ref{sec:setup}, \todo{harmonized re-run}); the full
model does not. \emph{Theirs:} DGFusion trains with ground-truth
log-depth supervision and CLIP condition meta-text, and CAFuser with
condition annotations; both use roughly a $2\times$ longer effective
training schedule (ours: 200 epochs). \emph{Both directions:} no
test-time augmentation on our side, matched input modalities, and
val-only checkpoint selection for every headline claim.
